\label{ch:methodology-and-system-overview}
This chapter presents the methodological framework and describes the design 
and development of MLguide, a web-based decision-support system for machine 
learning method selection. 
Section~\ref{sec:dsr} introduces Design Science Research as the overarching 
framework. 
Section~\ref{sec:system-overview} gives a high-level overview of MLguide 
and its main components. 
The following sections describe the design goals, the development process, 
and the user testing approach.

\section{Design Science Research}
\label{sec:dsr}

This thesis uses Design Science Research (DSR) as its methodological framework, 
following \textcite{peffers2007design_science_research_methodology} and 
\textcite{hevner2004design_science_information_systems_research}. 
DSR is a well-established approach for developing and evaluating decision support 
systems \parencite{arnott2012design_science_decision_support_systems_research}, 
and has been applied in recent work combining machine learning with user-facing 
decision support 
\parencite{landwehr2022design_knowledge_deep_learning, dellermann2018design_principles_hybrid_intelligence}. 
In DSR, knowledge is gained by building and evaluating an artifact, 
producing both a usable system and knowledge that can inform future work of a similar kind
\parencite{landwehr2022design_knowledge_deep_learning, hevner2004design_science_information_systems_research}. 
The DSR approach is therefore suited for this work, where the research questions are investigated 
through the design and evaluation of the artifact, MLguide.

The DSR process consists of six activities: problem identification, defining objectives, 
design and development, demonstration, evaluation, and communication 
\parencite{peffers2007design_science_research_methodology}. 
These activities are nominally sequential, but iteration is expected. 
After evaluation, the researcher may return to design and development to improve 
the artifact before proceeding \parencite{peffers2007design_science_research_methodology}. 
Figure~\ref{fig:dsr-methodology-overview} shows how these activities map to the structure 
of this thesis.

\begin{figure}[H]
    \centering
    \includegraphics[width=\textwidth]{Figures/03-research-methodology-system-overview/dsrm-mlguide.pdf}
    \caption[DSR methodology]{Design Science Research Methodology adapted from \textcite{peffers2007design_science_research_methodology}.}
    \label{fig:dsr-methodology-overview}
\end{figure}

The process in this thesis followed a design-centered initiation, where
development of the system started before the research questions were fully formalised
\parencite{peffers2007design_science_research_methodology}.

\section{System Overview}
\label{sec:system-overview}

MLguide is a web-based decision-support system designed to help users select 
machine learning methods for a given problem. 
The user starts by describing their problem with text input,
and can provide filters to narrow the scope of the problem, such as ML area (earlier referred to as paradigm in this thesis) and task type. 
Figure~\ref{fig:mlguide-overview} gives a high-level overview of MLguide and its main components. 

\begin{figure}[H]
    \centering
    \includegraphics[width=\textwidth]{Figures/03-research-methodology-system-overview/mlguide-system-overview.pdf}
    \caption[MLguide system overview]{System Overview of MLguide.}
    \label{fig:mlguide-overview}
\end{figure}

The system uses two knowledge sources. 
The first is a structured ML knowledge graph containing information about machine learning methods and tasks, 
including links between methods and the research articles they appear in, among other things. 
The second is a database of research articles from the literature. 
With the help of these knowledge sources, the recommendation engine produces three outputs: 
a ranked list of ML methods, a set of supporting articles, 
and a starter notebook, a Jupyter notebook file with code generated by MLguide as a starting point for implementation.

\section{Design Goals and Motivation}
\label{sec:design-goals-and-motivation}

MLguide was developed as a direct continuation of the work done in the pre-master project \parencite{bergsto2026mlselection}.
The developed roadmap outlines a process for filtering and ranking 
research literature based on user-defined context parameters, including 
ML area, product lifecycle (PLC) phase, and application area. 
This is illustrated with a simplified version of the roadmap in Figure~\ref{fig:ml-roadmap-pre-master}.
In the pre-master project, the main focus was on collecting and analysing 
research literature. Several directions for future development were identified, 
including better support for problem formulation, ontology-based ML knowledge 
representation, support for translating recommendations into an 
ML implementation, and incorporating user feedback to improve method 
recommendations over time \parencite{bergsto2026mlselection}.

\begin{figure}[H]
    \centering
    \includegraphics[width=0.4\textwidth]{Figures/04-system-design-development/roadmap-pre-master.pdf}
    \caption[ML selection roadmap]{Simplified ML selection roadmap from the pre-master project \parencite{bergsto2026mlselection}.}
    \label{fig:ml-roadmap-pre-master}
\end{figure}

MLguide addresses these directions by extending the original roadmap into
a web-based decision-support system. The main design goals are to support 
users in formulating their problem, to recommend relevant ML methods based 
on both structured knowledge and research literature, to provide a 
starting point for implementation through generated notebooks, and to 
incorporate user feedback on recommendations. 
The system is primarily aimed at engineers and students who are new to 
machine learning and need structured guidance in choosing a suitable method.
A web-based interface was chosen to 
lower the threshold for use, making it easy for other users to access and test
the system without installation.

\section{Incremental and Iterative Development}
\label{sec:incremental-and-iterative-development}

MLguide was developed incrementally over the course of the master's project. 
New features and components were built and deployed one at a time, each adding 
to the working system. This approach is consistent with incremental software 
development, where the system grows through a series of versions rather than 
being built all at once \parencite[33]{sommerville2011software_engineering}.

Figure~\ref{fig:incremental-development} illustrates this process. 
Planning, development, and self-evaluation were carried out as concurrent 
activities, producing successive versions of the system. 
Throughout development, individual features and the overall system output 
were continuously tested to check that behaviour was correct and results were useful.

The development also had an iterative character. Feedback from the thesis 
supervisor informed adjustments to the system design and priorities throughout 
the development period. 
In addition, two rounds of user testing were conducted, 
and the findings were used to evaluate and adjust the system. 
This iterative cycle is consistent with the DSR process described by 
\textcite{peffers2007design_science_research_methodology}, where evaluation 
results inform a return to design and development. The user testing setup is described in the following section.

\begin{figure}[H]
    \centering
    \includegraphics[width=0.8\textwidth]{Figures/04-system-design-development/incremental-development.pdf}
    \caption[Incremental development process]{Illustration of the incremental development process, based on
    \textcite[33]{sommerville2011software_engineering}.}
    \label{fig:incremental-development}
\end{figure}

\section{User Testing}
\label{sec:user-testing}

User testing was conducted in collaboration with the course TMM4128 Machine
Learning for Engineers at the Norwegian University of Science and Technology
(NTNU). The course provided a natural opportunity to test the system with
its intended user group. Students used MLguide as part of their course
projects. The participants were engineering students with limited prior
experience with machine learning. Feedback could be collected through a feedback form
accessible in MLguide and from the students' project reports.
The two course projects are summarised in Table~\ref{tab:user-testing}.

\begin{table}[H]
\centering
\caption[User testing sessions]{Description of the two user testing sessions conducted in collaboration with the course TMM4128 Machine Learning for Engineers at NTNU.}
\label{tab:user-testing}
\begin{tabular}{l p{5.5cm} p{5.5cm}}
\toprule
 & \textbf{Project 1} & \textbf{Project 2} \\
\midrule
\textbf{Topic} & Anomaly detection in time series data & Reinforcement learning in product lifecycle \\
\addlinespace
\textbf{ML area} & Unsupervised learning & Reinforcement learning \\
\addlinespace
\textbf{Task} & Detect anomalies using unsupervised learning techniques and compare performance & Analyse RL applications across lifecycle phases and implement a sustainability agent \\
\addlinespace
\textbf{MLguide task} & Test the system on their problem and report on relevance and usefulness of results & Assess the applicability of MLguide to their problem \\
\addlinespace
\textbf{Deadline} & March 2026 & April 2026 \\
\bottomrule
\end{tabular}
\end{table}

The two projects differed in scope and focus. The first project had a
well-defined problem: detecting anomalies in time series data using
unsupervised learning. Students were asked to compare at least three
techniques, which made MLguide a natural fit. The system could be used
early in the project to get method recommendations, or later as a
comparison against methods already selected.

The second project was more open-ended. Students were asked to analyse
applications of reinforcement learning across four product lifecycle
phases: product design, manufacturing, product use, and recycling,
and to implement a simulated RL environment and a sustainability agent.
MLguide could be used in two ways in this context: to find relevant
articles by submitting problem descriptions related to different
lifecycle phases, and to get algorithm recommendations for the
implementation part of the project.

A simple overview of the development timeline, including the user testing sessions, is illustrated in Figure~\ref{fig:development-timeline}.

\begin{figure}[H]
    \centering
    \includegraphics[width=0.8\textwidth]{Figures/04-system-design-development/development_timeline.pdf}
    \caption[Development timeline]{Development timeline of MLguide.}
    \label{fig:development-timeline}
\end{figure}
